% =====================================================================
%  Semestrální práce – předmět Hlasové dialogové systémy (HDS)
%  Hlasový kuchyňský asistent
%
%  Překlad:  pdflatex dokumentace.tex   (2x kvůli obsahu a odkazům)
% =====================================================================
\documentclass[11pt,a4paper]{article}

\usepackage[T1]{fontenc}
\usepackage[utf8]{inputenc}
\usepackage[czech]{babel}
\usepackage{lmodern}
\usepackage[a4paper,margin=2.5cm]{geometry}
\usepackage{graphicx}
\usepackage{xcolor}
\usepackage{booktabs}
\usepackage{enumitem}
\usepackage{listings}
\usepackage{tikz}
\usetikzlibrary{shapes.geometric,arrows.meta,positioning,calc}
\usepackage[hidelinks]{hyperref}
\usepackage{fancyhdr}

% --- Barvy ----------------------------------------------------------
\definecolor{accent}{RGB}{232,93,42}
\definecolor{codebg}{RGB}{248,247,243}
\definecolor{codekw}{RGB}{0,92,197}
\definecolor{codecom}{RGB}{106,115,89}
\definecolor{codestr}{RGB}{163,21,69}

% --- Nastavení výpisů kódu (vč. českých znaků v komentářích) ---------
\lstset{
  backgroundcolor=\color{codebg},
  basicstyle=\ttfamily\footnotesize,
  keywordstyle=\color{codekw}\bfseries,
  commentstyle=\color{codecom}\itshape,
  stringstyle=\color{codestr},
  numberstyle=\tiny\color{gray},
  numbers=left,
  numbersep=8pt,
  frame=single,
  rulecolor=\color{gray!40},
  breaklines=true,
  showstringspaces=false,
  tabsize=4,
  captionpos=b,
  inputencoding=utf8,
  extendedchars=true,
  literate=%
    {á}{{\'a}}1 {é}{{\'e}}1 {í}{{\'i}}1 {ó}{{\'o}}1 {ú}{{\'u}}1 {ý}{{\'y}}1
    {Á}{{\'A}}1 {É}{{\'E}}1 {Í}{{\'I}}1 {Ó}{{\'O}}1 {Ú}{{\'U}}1 {Ý}{{\'Y}}1
    {č}{{\v c}}1 {ď}{{\v d}}1 {ě}{{\v e}}1 {ň}{{\v n}}1 {ř}{{\v r}}1
    {š}{{\v s}}1 {ť}{{\v t}}1 {ž}{{\v z}}1 {ů}{{\r u}}1
    {Č}{{\v C}}1 {Ď}{{\v D}}1 {Ě}{{\v E}}1 {Ň}{{\v N}}1 {Ř}{{\v R}}1
    {Š}{{\v S}}1 {Ť}{{\v T}}1 {Ž}{{\v Z}}1 {Ů}{{\r U}}1
}

% --- Záhlaví/zápatí --------------------------------------------------
\pagestyle{fancy}
\fancyhf{}
\lhead{\small Hlasový kuchyňský asistent}
\rhead{\small Hlasové dialogové systémy}
\cfoot{\thepage}
\renewcommand{\headrulewidth}{0.4pt}

\begin{document}

% =====================================================================
%  Titulní strana
% =====================================================================
\begin{titlepage}
\centering
\vspace*{2cm}
{\Large\scshape Západočeská univerzita v Plzni\par}
{\large Fakulta aplikovaných věd\par}
\vspace{0.4cm}
{\large Předmět: Hlasové dialogové systémy (HDS)\par}
\vfill
{\color{accent}\rule{\linewidth}{1.4pt}}\par
\vspace{0.6cm}
{\Huge\bfseries Hlasový kuchyňský asistent\par}
\vspace{0.4cm}
{\large Hlasový dialogový systém pro provádění uživatele receptem\\ na bázi LLM (Ollama), RAG a SpeechCloud\par}
\vspace{0.6cm}
{\color{accent}\rule{\linewidth}{1.4pt}}\par
\vfill
{\large\bfseries Semestrální práce\par}
\vspace{2cm}
\begin{flushright}
\large
\textbf{Autor:} David Král\\
\textbf{E-mail:} davidkral03@gmail.com\\
\textbf{Akademický rok:} 2025/2026
\end{flushright}
\vspace{1cm}
{\large \today\par}
\end{titlepage}

% =====================================================================
%  Anotace + obsah
% =====================================================================
\begin{abstract}
\noindent
Tato práce popisuje návrh a implementaci \textbf{hlasového kuchyňského asistenta} --
webové aplikace, která uživatele interaktivně provede vybraným receptem krok za
krokem. Uživatel nahraje libovolný recept v textové podobě; systém ho automaticky
rozloží na suroviny a jednotlivé kroky, uloží do vyhledávací (RAG) databáze a poté
ho receptem provází. Aplikace umožňuje ovládání jak \emph{psaním} v chatu, tak
\emph{hlasem} -- s rozpoznáváním řeči (ASR) a syntézou řeči (TTS) zajišťovanými
platformou SpeechCloud. Porozumění a generování odpovědí obstarává velký jazykový
model provozovaný přes Ollamu. Klíčovou vlastností je \textbf{deterministické sledování
postupu}: aplikace si v kódu drží, ve kterém kroku se uživatel nachází, takže se ve
vaření nikdy neztratí. Uživatelské rozhraní je inspirováno moderními konverzačními
asistenty -- nabízí klasický chat s možností přepnutí do plnohodnotného hlasového
režimu s přehledem konverzace.
\end{abstract}

\tableofcontents
\newpage

% =====================================================================
\section{Úvod}
% =====================================================================
Vaření podle receptu na obrazovce má praktická úskalí: ruce má člověk zaměstnané
a často špinavé, displej zhasíná a v dlouhém textu se snadno ztratí orientace,
ve kterém kroku se právě nachází. Hlasové dialogové systémy nabízejí přirozené
řešení -- uživatel může s aplikací mluvit a poslouchat ji, aniž by se musel
dotýkat zařízení.

Cílem této semestrální práce bylo navrhnout a implementovat hlasového asistenta,
který:
\begin{itemize}[noitemsep]
  \item přijme od uživatele libovolný recept (vložený text nebo soubor),
  \item automaticky z něj vytěží suroviny a posloupnost kroků,
  \item uloží recept do RAG databáze a umí odpovídat na dotazy k němu,
  \item provází uživatele krokem za krokem a \emph{nikdy se v krocích neztratí},
  \item umí poslouchat (ASR) i mluvit (TTS),
  \item nabízí příjemné moderní rozhraní s chatem i hlasovým režimem.
\end{itemize}

Práce navazuje na příklady probírané v předmětu (framework SpeechCloud pro řízení
dialogu, připojení k Ollamě a ukázka RAG) a spojuje je do jedné funkční aplikace.

% =====================================================================
\section{Zadání a požadavky}
% =====================================================================
Funkční požadavky na systém lze shrnout takto:
\begin{enumerate}[noitemsep]
  \item \textbf{Nahrání receptu} -- uživatel vloží text receptu (suroviny i postup).
  \item \textbf{Strukturalizace} -- systém recept rozdělí na suroviny a kroky.
  \item \textbf{RAG} -- recept se uloží do vyhledávací databáze pro zodpovídání dotazů.
  \item \textbf{Provázení krok za krokem} -- s udržením aktuální pozice v postupu.
  \item \textbf{Dva režimy ovládání} -- textový chat a hlasové ovládání.
  \item \textbf{Hlasový vstup/výstup} -- rozpoznávání a syntéza řeči.
  \item \textbf{Robustnost} -- model se nesmí v krocích ztratit ani přeskakovat.
\end{enumerate}

Nefunkční požadavky: srozumitelné a vizuálně příjemné rozhraní, odezva v reálném
čase, odpovědi v češtině a snadné spuštění.

% =====================================================================
\section{Teoretický základ}
% =====================================================================

\subsection{Hlasové dialogové systémy}
Hlasový dialogový systém (HDS) je systém, který s uživatelem komunikuje v přirozeném
jazyce pomocí řeči. Typicky se skládá z rozpoznávače řeči (ASR -- Automatic Speech
Recognition), porozumění (NLU/SLU), řízení dialogu (Dialog Management), generování
odpovědi (NLG) a syntézy řeči (TTS -- Text-to-Speech). V této práci roli NLU i NLG
přebírá velký jazykový model, řízení dialogu zajišťuje vlastní logika v kombinaci
s LLM a ASR/TTS poskytuje SpeechCloud.

\subsection{Rozpoznávání a syntéza řeči -- SpeechCloud}
\textbf{SpeechCloud} je platforma Katedry kybernetiky ZČU poskytující ASR a TTS jako
službu. V prohlížeči běží klient (\texttt{speechcloud-3.0.js}), který přes WebRTC/SIP
zachytává zvuk z mikrofonu a přehrává syntetizovanou řeč. Serverová část dialogu se
SpeechCloudem komunikuje přes WebSocket pomocí frameworku \texttt{dialog.py}, který
nabízí pohodlné metody jako \texttt{synthesize\_and\_wait()} (řekni a počkej) nebo
\texttt{recognize\_and\_wait\_for\_asr\_result()} (poslouchej a vrať rozpoznaný text).

\subsection{Velké jazykové modely a Ollama}
\textbf{Ollama} je server pro provoz velkých jazykových modelů (LLM). V práci je využit
model \texttt{gemma3:12b} pro porozumění a generování odpovědí v češtině a dedikovaný
model \texttt{qwen3-embedding} pro výpočet embeddingů. Důležitou schopností je
\emph{strukturovaný výstup}: modelu lze předat JSON schéma a vynutit, aby odpověď
přesně odpovídala dané struktuře. Toho práce využívá pro spolehlivé parsování receptu
i pro klasifikaci záměru uživatele.

\subsection{Retrieval-Augmented Generation (RAG)}
RAG kombinuje vyhledávání v dokumentech s generováním odpovědí. Dokument (zde recept)
se rozdělí na úseky, pro každý se spočítá \emph{embedding} (číselný vektor zachycující
význam) a uloží do databáze. Při dotazu se spočítá embedding dotazu a kosinovou
podobností se najdou nejrelevantnější úseky, které se předají LLM jako kontext.
Model tak odpovídá na základě konkrétního receptu, nikoli jen z obecné znalosti.

\subsection{Řízení dialogu a sledování stavu}
Spoléhat při krokování receptu výhradně na LLM je rizikové -- model může krok
přeskočit, zopakovat nebo „zapomenout``, kde se nachází. Práce proto odděluje
\emph{rozhodnutí o záměru} (co uživatel chce) od \emph{provedení navigace}. Záměr
určuje LLM (případně rychlá pravidla), ale samotný posun v krocích je čistě
deterministická operace nad explicitně drženým stavem. Tím je vyloučeno, že by se
asistent v postupu ztratil.

% =====================================================================
\section{Architektura systému}
% =====================================================================
Systém má dvě hlavní části: \textbf{frontend} (běží v prohlížeči) a \textbf{backend}
(server v Pythonu nad frameworkem Tornado). Frontend přes SpeechCloud zajišťuje
mikrofon a reproduktor a přes WebSocket komunikuje s dialogovým manažerem na serveru.
Backend řídí dialog, volá Ollamu (LLM + embeddingy) a spravuje RAG i stav postupu.

\begin{figure}[h]
\centering
\begin{tikzpicture}[
  font=\small,
  box/.style={draw,rounded corners,align=center,minimum height=1cm,inner sep=5pt,fill=white},
  acc/.style={draw=accent,line width=1pt},
  arr/.style={-{Latex[length=2.5mm]},thick}
]
\node[box,acc] (chat) {Chat + Hlasový režim\\\scriptsize (HTML/CSS/JS)};
\node[box,below=0.7cm of chat] (sc) {SpeechCloud klient\\\scriptsize mikrofon / reproduktor};
\node[box,acc,right=3.2cm of chat] (dm) {Dialogový manažer\\\scriptsize KitchenDialog (Tornado)};
\node[box,below=0.7cm of dm] (state) {RecipeState\\\scriptsize sledování kroku};
\node[box,right=2.6cm of dm] (llm) {Ollama\\\scriptsize LLM + embeddingy};
\node[box,below=0.7cm of llm] (rag) {RecipeRAG\\\scriptsize vektorové vyhledávání};

\draw[arr] (chat) -- node[above]{\scriptsize WS /ws} (dm);
\draw[arr] (dm) -- (chat);
\draw[arr] (sc) -- node[below,sloped]{\scriptsize ASR/TTS} ($(dm.south west)+(0,-0.1)$);
\draw[arr] (dm) -- (llm);
\draw[arr] (llm) -- (dm);
\draw[arr] (dm) -- (state);
\draw[arr] (dm) -- (rag);
\draw[arr] (rag) -- (llm);
\end{tikzpicture}
\caption{Zjednodušené schéma architektury systému.}
\label{fig:arch}
\end{figure}

\paragraph{Komunikace frontend $\leftrightarrow$ backend.}
Probíhá obecným kanálem SpeechCloud. Frontend posílá data zprávou
\texttt{dm\_send\_message}, backend odpovídá zprávou \texttt{dm\_receive\_message}.
Nad tímto kanálem je definován jednoduchý aplikační protokol (tabulka~\ref{tab:proto}).

\begin{table}[h]
\centering
\small
\begin{tabular}{@{}lll@{}}
\toprule
\textbf{Směr} & \textbf{Typ zprávy} & \textbf{Význam} \\
\midrule
FE\,$\to$\,BE & \texttt{set\_recipe}      & nahrání textu receptu \\
FE\,$\to$\,BE & \texttt{user\_text}       & napsaná zpráva uživatele \\
FE\,$\to$\,BE & \texttt{set\_voice\_mode} & zapnutí/vypnutí hlasového režimu \\
FE\,$\to$\,BE & \texttt{reset}            & vymazání konverzace \\
BE\,$\to$\,FE & \texttt{assistant}        & odpověď asistenta \\
BE\,$\to$\,FE & \texttt{user\_speech}     & přepis rozpoznané řeči \\
BE\,$\to$\,FE & \texttt{recipe\_loaded}   & strukturovaný recept \\
BE\,$\to$\,FE & \texttt{progress}         & aktuální krok / počet kroků \\
BE\,$\to$\,FE & \texttt{status}           & stav (přemýšlím/mluvím/poslouchám) \\
\bottomrule
\end{tabular}
\caption{Aplikační protokol mezi frontendem (FE) a backendem (BE).}
\label{tab:proto}
\end{table}

% =====================================================================
\section{Implementace}
% =====================================================================
Projekt je členěn do několika souborů:
\begin{itemize}[noitemsep]
  \item \texttt{app.py} -- dialogový manažer \texttt{KitchenDialog} a spuštění serveru,
  \item \texttt{recipe.py} -- parsování receptu, RAG a sledování kroků,
  \item \texttt{config.py} -- konfigurace (server, model, hlas, port),
  \item \texttt{dialog.py} -- framework SpeechCloud (z materiálů předmětu),
  \item \texttt{\_event\_emitter.py} -- záložní EventEmitter,
  \item \texttt{static/} -- frontend (\texttt{index.html}, \texttt{style.css}, \texttt{app.js}).
\end{itemize}

\subsection{Dialogový manažer a souběžnost}
Jádrem backendu je třída \texttt{KitchenDialog} dědící z \texttt{Dialog}. Její hlavní
smyčka musí současně reagovat na dvě události: příchod zprávy z webu (text, příkaz)
\emph{a} rozpoznání řeči. Řeší to souběžným čekáním na obě úlohy pomocí
\texttt{asyncio.wait} s návratem při první dokončené úloze. Rozpoznávání se spouští
jen v hlasovém režimu a jen tehdy, když asistent zrovna nemluví ani nepřemýšlí
(příznak \texttt{busy}), aby neslyšel sám sebe.

\begin{lstlisting}[language=Python,caption={Hlavní smyčka -- souběžné čekání na text i řeč.}]
self._msg_task = asyncio.ensure_future(self.pop_message())
while True:
    tasks = {self._msg_task}
    if self.voice_mode and not self.busy:
        if self._asr_task is None:
            self._asr_task = asyncio.ensure_future(
                self.recognize_and_wait_for_asr_result(timeout=config.ASR_TIMEOUT))
        tasks.add(self._asr_task)

    done, _ = await asyncio.wait(tasks, return_when=asyncio.FIRST_COMPLETED)

    if self._msg_task in done:                 # přišla zpráva z webu
        msg = self._msg_task.result()
        self._msg_task = asyncio.ensure_future(self.pop_message())
        await self.handle_message(msg)

    if self._asr_task is not None and self._asr_task in done:  # rozpoznaná řeč
        result = self._asr_task.result()
        self._asr_task = None
        await self.handle_voice_result(result)
\end{lstlisting}

\subsection{Parsování receptu (strukturovaný výstup)}
Recept se zpracovává jediným voláním LLM s vynuceným JSON schématem
(\texttt{title}, \texttt{ingredients}, \texttt{steps}). Systémový prompt modelu
ukládá rozdělit dlouhá souvětí na samostatné, jednotlivě proveditelné kroky -- to je
klíčové pro pozdější krokování. Pokud by volání LLM selhalo, použije se deterministický
záložní parser (dělení podle nadpisů „Suroviny`` / „Postup`` a číslování), takže
funkce vrátí použitelný výsledek vždy.

\begin{lstlisting}[language=Python,caption={JSON schéma pro strukturovaný výstup parsování.}]
RECIPE_SCHEMA = {
    "type": "object",
    "properties": {
        "title":       {"type": "string"},
        "ingredients": {"type": "array", "items": {"type": "string"}},
        "steps":       {"type": "array", "items": {"type": "string"}},
    },
    "required": ["title", "ingredients", "steps"],
}
\end{lstlisting}

\subsection{RAG nad receptem}
Třída \texttt{RecipeRAG} rozdělí recept na úseky (blok surovin a každý krok zvlášť),
pro každý spočítá embedding modelem \texttt{qwen3-embedding} a uloží jej. Při dotazu
se spočítá embedding dotazu a kosinovou podobností se vyberou nejrelevantnější úseky.
Když embeddingy nejsou dostupné, systém se elegantně přepne na vyhledávání podle
překryvu klíčových slov -- aplikace tedy funguje i bez embeddingové služby.

\begin{lstlisting}[language=Python,caption={Vyhledání relevantních úseků receptu.}]
def retrieve(self, query, k=3):
    if self.use_embeddings:
        qvec = self.client.embed(model=self.embed_model, input=[query])["embeddings"][0]
        scored = [(_cosine(qvec, v), c) for v, c in zip(self.vectors, self.chunks)]
        scored.sort(key=lambda x: x[0], reverse=True)
        return [c for _, c in scored[:k]]
    return self._keyword_retrieve(query, k)   # záloha bez embeddingů
\end{lstlisting}

\subsection{Deterministické sledování kroků}
Pozici v receptu drží třída \texttt{RecipeState} jako index do seznamu kroků. Navigace
(\texttt{start}, \texttt{next}, \texttt{previous}, \texttt{goto}) je ošetřena tak, aby
index nikdy nevypadl z platného rozsahu. Díky tomu je posun mezi kroky zcela
předvídatelný a nezávislý na „náladě`` jazykového modelu -- to je odpověď na požadavek,
aby se asistent v krocích neztratil.

\begin{lstlisting}[language=Python,caption={Bezpečná navigace nad stavem postupu.}]
def next(self):
    if not self.steps:
        return None
    if self.index < self.total - 1:   # nikdy nepřekročíme poslední krok
        self.index += 1
    return self.current_text()

def goto(self, number):               # 1-based číslo kroku od uživatele
    if not self.steps:
        return None
    self.index = max(0, min(self.total - 1, number - 1))
    return self.current_text()
\end{lstlisting}

\subsection{Klasifikace záměru}
Každá promluva uživatele se zařadí do jednoho ze záměrů: \texttt{start}, \texttt{next},
\texttt{previous}, \texttt{repeat}, \texttt{goto}, \texttt{ingredients}, \texttt{finish}
nebo \texttt{question}. Pro běžné krátké příkazy („dál``, „co dál``, „zopakuj`` $\dots$)
se použije \emph{rychlá cesta} bez LLM -- okamžitě a spolehlivě. Delší či nejednoznačné
promluvy klasifikuje LLM se strukturovaným výstupem. Záměr \texttt{question} spustí RAG
a vygenerování odpovědi; navigační záměry jen deterministicky posunou stav a přečtou
příslušný krok.

\begin{lstlisting}[language=Python,caption={Rychlá klasifikace běžných příkazů (zkráceno).}]
_QUICK_INTENTS = {
    "next":     ["další", "dál", "co dál", "pokračuj", "hotovo", ...],
    "previous": ["zpět", "předchozí", "vrať se", ...],
    "repeat":   ["zopakuj", "co teď", "znovu", ...],
    "start":    ["začni", "pojďme na to", "start", ...],
    ...
}
\end{lstlisting}

\subsection{Frontend}
Rozhraní je inspirováno konverzačními asistenty. Hlavní obrazovku tvoří chat
s polem pro psaní a postranní panel se strukturou receptu, v němž je zvýrazněn
aktuální krok a zobrazen ukazatel postupu. Tlačítkem mikrofonu lze přejít do
\emph{hlasového režimu} -- celoobrazovkového overlaye s animovaným indikátorem stavu
(poslouchám / přemýšlím / mluvím), živým přepisem řeči a přehledem konverzace.
O přehrávání syntézy a zachytávání mikrofonu se stará klient SpeechCloud; veškerá
další logika běží ve skriptu \texttt{app.js}.

% =====================================================================
\section{Použité technologie a konfigurace}
% =====================================================================
\begin{table}[h]
\centering\small
\begin{tabular}{@{}ll@{}}
\toprule
\textbf{Komponenta} & \textbf{Technologie} \\
\midrule
Webový server / WS    & Python, Tornado (\texttt{dialog.py}) \\
ASR + TTS             & SpeechCloud (\texttt{speechcloud-3.0.js}) \\
LLM                   & Ollama, model \texttt{gemma3:12b} \\
Embeddingy (RAG)      & Ollama, model \texttt{qwen3-embedding} \\
Frontend              & HTML5, CSS3, JavaScript, jQuery \\
\bottomrule
\end{tabular}
\caption{Přehled použitých technologií.}
\end{table}

Veškeré nastavení (adresa a přihlašovací údaje k Ollamě, název modelu, hlas pro TTS,
port serveru) je soustředěno v souboru \texttt{config.py}, takže se dá snadno změnit
na jednom místě.

% =====================================================================
\section{Spuštění a ovládání}
% =====================================================================
\subsection{Instalace a spuštění}
\begin{lstlisting}[caption={Instalace závislostí a spuštění serveru.}]
python3 -m venv venv
source venv/bin/activate
pip install -r requirements.txt
python app.py
# v prohlížeči otevřít:  http://localhost:8888/static/index.html
\end{lstlisting}
Hlasový režim vyžaduje mikrofon a kontext \texttt{localhost} (nebo HTTPS), jinak
prohlížeč přístup k mikrofonu zablokuje.

\subsection{Ovládání}
\begin{enumerate}[noitemsep]
  \item Tlačítkem \emph{Nahrát recept} se vloží text receptu (nebo ukázka).
  \item V chatu lze psát příkazy (\emph{začni}, \emph{další}, \emph{zpět},
        \emph{zopakuj}, \emph{krok tři}, \emph{suroviny}) i volné dotazy
        (\emph{kolik mouky?}, \emph{čím nahradit máslo?}).
  \item Ikonou mikrofonu se zapne hlasový režim -- stejné příkazy lze říkat.
\end{enumerate}

% =====================================================================
\section{Příklad interakce}
% =====================================================================
\begin{quote}\small
\textbf{Asistent:} Načetl jsem recept Palačinky. Má 11 kroků a 6 surovin. Až budete
chtít, řekněte „začni``.\\
\textbf{Uživatel:} Začni.\\
\textbf{Asistent:} Do mísy prosejte mouku a přidejte špetku soli a cukr.\\
\textbf{Uživatel:} Kolik mléka tam dám?\\
\textbf{Asistent:} Podle receptu celkem 500 ml; nejdřív přilijte zhruba polovinu.\\
\textbf{Uživatel:} Co dál?\\
\textbf{Asistent:} Rozklepněte vejce a přilijte zhruba polovinu mléka.
\end{quote}
Ukázka demonstruje, že dotaz mimo postup (množství mléka) nezmění pozici v receptu --
po „co dál`` asistent korektně pokračuje dalším krokem.

% =====================================================================
\section{Diskuse a možná rozšíření}
% =====================================================================
Oddělení rozhodnutí o záměru od provedení navigace se ukázalo jako zásadní pro
spolehlivost -- asistent se v krocích neztrácí ani při delších konverzacích. Mezi
omezení patří závislost na dostupnosti vzdálených služeb (SpeechCloud, Ollama) a
ukončení relace SpeechCloud při delší nečinnosti, které aplikace řeší nabídkou
opětovného připojení.

Možná rozšíření:
\begin{itemize}[noitemsep]
  \item časovače pro kroky s definovanou dobou (např. „peče se 20 minut``),
  \item přepočet množství surovin podle počtu porcí,
  \item podpora více receptů a jejich kombinování,
  \item perzistentní RAG databáze (např. FAISS) sdílená mezi relacemi.
\end{itemize}

% =====================================================================
\section{Závěr}
% =====================================================================
Cílem práce bylo vytvořit funkčního hlasového kuchyňského asistenta, který uživatele
spolehlivě provede receptem. Výsledná aplikace splňuje všechny stanovené požadavky:
přijme a strukturalizuje recept, uloží jej do RAG databáze, provází uživatele krok za
krokem s deterministickým udržením pozice a umožňuje ovládání psaním i hlasem v
příjemném moderním rozhraní. Práce zároveň ukázala, jak vhodně zkombinovat velký
jazykový model s jednoduchou, ale spolehlivou stavovou logikou tak, aby výsledek byl
robustní i prakticky použitelný.

% =====================================================================
\begin{thebibliography}{9}
\bibitem{speechcloud} Katedra kybernetiky ZČU. \emph{SpeechCloud -- platforma pro ASR a TTS}. Plzeň.
\bibitem{ollama} Ollama. \emph{Run large language models locally}. \url{https://ollama.com}.
\bibitem{rag} Lewis, P. et al. \emph{Retrieval-Augmented Generation for Knowledge-Intensive NLP Tasks}. 2020.
\bibitem{tornado} \emph{Tornado Web Framework}. \url{https://www.tornadoweb.org}.
\end{thebibliography}

\end{document}
